\subsubsection{Déclaration du type}
\type{CodeBinaire}{\chaine}

\subsubsection{obtenirCodeBinaire}

\begin{algorithme}
    \fonction{obtenirCodeBinaire}{a : Arbre Huffman, car : \caractere}{CodeBinaire}
    {non(estVide(a))}
    {}
    {\sialorssinon{estFeuille(a)}
    {\retourner('0')}
    \retourner{ObtenirCodeBinaireR(a,car)}}

\subsubsection{obtenirCodeBinaireR}

\begin{algorithme}
    \fonction{obtenirCodeBinaireR}{a : Arbre Huffman, car : \caractere}{CodeBinaire}
    {non(estVide(a))}
    {}
    {\sialorssinon{estFeuille(a)}
    {\retourner{''}}
    {\sialorssinon{estPresent(car, obtenirFilsGauche(a))}
    {\retourner{'0' + obtenirCodeBinaireR(obtenirFilsGauche(a),car)}}
    {\retourner{'1' + obtenirCodeBinaireR(obtenirFilsDroit(a),car)}}}}
\end{algorithme}

\subsubsection{estVide}

\begin{algorithme}
    \fonction{estVide}{code : CodeBinaire}{\booleen}
    {}
    {}
    {\retourner{code = ''}}
\end{algorithme}

\subsubsection{obtenirIemeBit}

\begin{algorithme}
    \fonction{obtenirIemeBit}{code : CodeBinaire, i : \naturelNonNul}{Bit}
    {i $>$ 0 et i \leq longueur(code)}
    {}
    {\retourner{code[i]}}
\end{algorithme}

\subsubsection{longueur}

\begin{algorithme}
    \fonction{longueur}{code : CodeBinaire}{\naturel}
    {}
    {l : naturel}
    {\affecter{l}{0}
    \pourChaque{car}{code}
    {\affecter{l}{l+1}}
    \retourner{l}}
\end{algorithme}